% ===================== UFO SHAPE TAXONOMY =====================
% Seventy-two reported shapes in strict A--Z order, nine to a page over eight
% pages, on the same
% card system as the Alien and Cryptid rosters but three columns wide, and
% over the Earth plate rather than the starfield or the swamp. The last line
% of every card is the useful one: what the shape nearly always turns out to
% be once somebody does the work.

The shape catalogue below is the part of this chapter I expect to be argued
with, so let me say plainly what it is and what it is not. It is not a
zoology of craft. It is an index of \emph{reported outlines} --- what witnesses
drew, what investigators filed, and what each outline has historically
resolved into. Read it as a lookup table. When a report lands on your desk,
find the shape, read the last line of the card, and eliminate that first. If
the ordinary explanation survives the elimination, you have something worth
your evening.

\earthfieldon
\rosterwidethree

\begin{multicols}{3}[\rosterhead{Classification of Common UFO Shapes}]

\ufoshapecard{\gacorn}{Acorn}{Discoid; Kecksburg, 1965}{A rounded body with a
collar or banded base. Kecksburg is the case that carries the type, and it is
also the case that shows how fast a shape hardens: the drawings became more
consistent \emph{after} the story was retold, not before.}{A re-entering space
object, most plausibly a spent probe or booster section.}

\ufoshapecard{\gamoeba}{Amorphous}{Formless; the residual category}{No stable
outline: reported as shifting, folding, or changing between glances. Every
classification system needs a bin for the reports that will not sit still, and
this is mine. It is also the least useful bin.}{Cloud, smoke, atmospheric
refraction, or the limits of night vision.}

\ufoshapecard{\gbarrel}{Barrel}{Cylindrical; industrial-area sightings}{A short
banded cylinder, sometimes reported as ribbed or seamed. Reports cluster near
ports, refineries and rail yards --- that is, near the places barrels
are.}{Industrial plant, a storage tank, or a shipping container under a crane.}

\ufoshapecard{\gbell}{Bell}{Discoid; various, 1940s onward}{Flared at the base,
narrowing to a crown. It appears in wartime and post-war European reports and
later in the exotic-engine literature, which is where the trouble starts: the
shape is now inseparable from claims about propulsion that nobody has
demonstrated.}{An illuminated tethered balloon, or an aircraft nose-on at
night.}

\ufoshapecard{\gblacktri}{Black Triangle}{Angular; Phoenix Lights, 1997}{The
night variant: unlit hull inferred from the pattern of lamps against the sky.
Phoenix showed exactly how this class fails --- one event was flares, another
was aircraft, and the \emph{shape} was supplied by the witnesses.}{Flares on
parachutes, an aircraft formation, or occlusion of stars misread as a hull.}

\ufoshapecard{\gboomerang}{Boomerang}{Angular; Hudson Valley, 1983--84}{A
V-shaped outline with a shallow apex, reported as gliding without sound. The
Hudson Valley wave was eventually traced in large part to a group of pilots
flying in formation for their own amusement, and they said so.}{Ultralight
formation flying, or migrating birds lit from below by city glow.}

\ufoshapecard{\gfractal}{Branching Form}{Formless; very rare}{A repeating
branched structure, reported chiefly in altered-state and abduction accounts
rather than in field sightings. The reporting context is the finding.}{Entoptic
phenomena --- form constants generated by the visual system itself.}

\ufoshapecard{\gbutterfly}{Butterfly}{Organic; sparse}{Paired lobed wings,
sometimes reported with markings. It is the class most likely to be an insect
near the lens, and the least likely to be anything at altitude.}{An insect
close to the camera, or a bird photographed mid-wingbeat.}

\ufoshapecard{\gcapsule}{Capsule}{Cylindrical; re-entry and test
windows}{Rounded at both ends, reported as descending steadily and sometimes as
under a canopy. Where a canopy is mentioned the case usually closes
immediately.}{A returning crew or cargo capsule, or a high-altitude test
article under parachute.}

\ufoshapecard{\gchain}{Chain of Lights}{Formation; Starlink era, post-2019}{A
regular line of evenly spaced lights holding station relative to each other.
This class exploded after 2019 and the cause is now so well established that
it has become a useful test of a reporting system's competence.}{A satellite
constellation train shortly after launch.}

\ufoshapecard{\gcigar}{Cigar}{Cylindrical; the 1896--97 airship flap}{An
elongated body, blunt at both ends, reported as huge and slow. It dominates
the pre-aviation record because it was the plausible future craft of the day.
Shape reports track the technology the public expects.}{An airship or
dirigible, a contrail lit end-on, or an aircraft seen along its axis.}

\ufoshapecard{\gdisc}{Classic Disc}{Discoid; Kenneth Arnold, 1947}{The
foundation shape of the whole subject, and an accident of language: Arnold
described the \emph{motion} as a saucer skipping on water, and the press gave
the world a saucer-shaped object. Every drawing filed afterwards was already
contaminated by that headline.}{A lenticular cloud, a weather balloon at
altitude, or the witness reproducing a picture they have already seen.}

\ufoshapecard{\gcloudcigar}{Cloud Cigar}{Composite; French wave, 1954}{An
elongated form with a cloudy, ill-defined envelope, sometimes reported
emitting smaller objects. Aimé Michel built a whole theory on these, and the
theory has aged considerably worse than the observations.}{A contrail or
lenticular formation, with aircraft nearby read as ejected craft.}

\ufoshapecard{\gcluster}{Cluster}{Formation; wave reports}{Several bodies
moving in loose company, described collectively as one phenomenon. Clusters
are where witness counts inflate fastest, and where the number reported is
least likely to be the number present.}{Balloon releases, sky lanterns at a
festival, or a drone display.}

\ufoshapecard{\gcomet}{Comet Form}{Luminous; spiral and plume events}{A bright
nucleus with a broad persistent tail, sometimes reported as spiralling. The
Norwegian spiral of 2009 was the making of this class and it was a failed
missile stage, confirmed and admitted.}{A rocket-stage propellant dump, or a
spiralling failed launch seen at range.}

\ufoshapecard{\gcone}{Cone}{Angular; nose-on and re-entry reports}{A circular
base drawn to an apex, reported both in flight and on the ground. It is what a
great many other shapes look like when they are pointed at you, so treat cone
reports as aspect-angle problems.}{A re-entry capsule or nose fairing, or an
aircraft seen nose-on.}

\ufoshapecard{\gcoolie}{Coolie Hat}{Discoid; Papua and Java reports, 1950s}{A
shallow cone over a flat base, reported chiefly from South-East Asia and the
Pacific in the early years. The regional concentration tells you it is at
least partly cultural: the witness reaches for a familiar silhouette.}{Cloud
form, or an aircraft seen from below with wing and fuselage merged.}

\ufoshapecard{\gcrescent}{Crescent}{Curved; wartime and Nazi-craft claims}{A
swept arc with concave trailing edge. It attracts the worst literature in the
subject, and the shape's own history --- a real German programme of flying
wings --- is more interesting than anything written about it since.}{A flying-wing
aircraft, or the moon's crescent through broken cloud.}

\ufoshapecard{\gcross}{Cross}{Angular; European reports, 1960s--70s}{Two
intersecting bars of light, reported as maintaining orientation. The reported
rigidity is the diagnostic feature and also the least reliable, because
witnesses are poor judges of rotation at distance.}{An aircraft head-on with
landing lights and wingtip strobes.}

\ufoshapecard{\gcube}{Cube in Sphere}{Geometric; contemporary, post-2015}{A
cubic form apparently inside a transparent shell, a description that arrived
late and spread fast online. Note the date. New shapes still enter the
catalogue, and they enter through the internet, not the sky.}{Composite or
edited imagery; the class has essentially no pre-digital record.}

\ufoshapecard{\gvoid}{Dark Mass}{Formless; nocturnal reports}{An area of sky
blacker than the sky, defined only by the stars it blocks. This is inference,
not observation, and it is the single easiest class to generate by accident
with a passing cloud.}{Cloud occluding starfield, or a large unlit balloon.}

\ufoshapecard{\gdelta}{Delta}{Angular; military airspace reports}{A swept
triangular planform with a notched trailing edge --- that is to say, a delta
wing. When a shape class matches a production aircraft this exactly, the
burden of proof sits with whoever claims it is not one.}{A delta-winged
aircraft; the type has been in service since the 1950s.}

\ufoshapecard{\gdiamond}{Diamond}{Geometric; Rendlesham Forest, 1980}{Two
opposed pyramids, reported as a lit lozenge low among trees. Rendlesham is the
case I would most like to be real and the case where the lighthouse timing
matches the testimony best.}{A lighthouse beam through trees, a bright star
low on the horizon, or re-entry debris.}

\ufoshapecard{\gdome}{Dome}{Discoid; ground-contact reports}{A hemisphere,
usually reported on or near the ground rather than in flight. It is the
landed-object shape, and it has the highest rate of physical-trace claims
attached to it, which at least makes it testable.}{Farm equipment, a grain
silo top, a tent, or a radome seen at distance.}

\ufoshapecard{\gdomedisc}{Domed Disc}{Discoid; Rex Heflin, 1965}{A disc with a
raised cupola, usually reported as glazed or windowed. It is the most drawn
variant in the archives because it is the most cinematic, and because a dome
gives the witness somewhere to put the occupants they have not seen.}{Hubcaps,
lens flare on a bright object, and outright models suspended on thread.}

\ufoshapecard{\gdumbbell}{Dumbbell}{Paired; scattered, all decades}{Two masses
joined by a slimmer waist. Whenever a report describes two connected objects,
ask first whether the connection was seen or inferred; it is nearly always
inferred.}{Two objects at different distances aligned by chance, or a tethered
pair of balloons.}

\ufoshapecard{\gegg}{Egg}{Ovoid; South American reports, 1950s--70s}{An ovoid
with one broader end, often reported as landed with occupants nearby. The
occupant claims come attached to the shape rather than the reverse, which
suggests the shape is doing narrative work.}{A hot-air balloon envelope
partially inflated, or a tank or vessel seen in a field.}

\ufoshapecard{\gfincyl}{Finned Cylinder}{Cylindrical; Cold War intrusion
reports}{A tube with tail surfaces, reported chiefly near military ranges and
coastlines. It is the shape most likely to have a boring answer, because the
places it is reported from are the places such objects are actually
launched.}{A test missile, sounding rocket, or target drone.}

\ufoshapecard{\gfireball}{Fireball}{Luminous; ubiquitous}{A brilliant
moving light with a visible head and short tail, usually reported for seconds
rather than minutes. It is the highest-volume report class worldwide and the
one with the highest resolution rate, which should tell you
something.}{A bolide meteor, or re-entering satellite debris burning up.}

\ufoshapecard{\ggyro}{Gyroscopic}{Composite; sparse}{Nested rotating rings
around a central body, usually reported by a single witness at length. Detail
of this order from one observer and none from anybody else is a reliable
warning sign.}{Confabulation, or a defocused light with rotational camera
movement.}

\ufoshapecard{\gnut}{Hexagonal Nut}{Geometric; sparse, modern}{A six-sided
plate with a central aperture, reported as flat and edge-lit. It appears
almost only in still photographs, never in multi-witness sightings, and that
asymmetry tells you where to look for the cause.}{A lens or sensor artefact
shaped by a hexagonal aperture blade set.}

\ufoshapecard{\ghourglass}{Hourglass}{Paired; sparse, modern}{Two cones meeting
at a narrow waist, sometimes reported as pulsing in step. Rare enough that
every entry in the class can be examined individually, and most of them are
photographs rather than sightings.}{An optical artefact --- an out-of-focus
light source with an iris or coma aberration.}

\ufoshapecard{\gicecream}{Ice Cream Cone}{Composite; ground-trace cases}{A
sphere seated in a cone, most often reported as landed on legs or a tripod.
The landing-gear detail is doing a lot of work here, and it is the detail
witnesses most often add after the fact.}{Farm or survey equipment, a
floodlight on a stand, or a hoax construction.}

\ufoshapecard{\gbowl}{Inverted Bowl}{Discoid; ground-level reports}{A shallow
concave form, reported open-side-down and usually stationary. It is the
second of the two landed-object shapes and it shares the first's habit of
turning out to be agricultural.}{A radome, a satellite dish, or a
polytunnel end seen at dusk.}

\ufoshapecard{\gjelly}{Jellyfish}{Organic; Iraq thermal footage, 2018}{A domed
body with trailing filaments. The declassified thermal clip revived this
class, and the honest answer is that a thermal sensor renders a mundane object
in an unfamiliar way and the eye supplies tentacles.}{Thermal blooming around
a balloon or drone, with sensor smear read as trailing limbs.}

\ufoshapecard{\gkite}{Kite Form}{Angular; daylight reports}{A four-sided
symmetrical outline, tapering below, reported as hovering and yawing. There is
a reason the shape shares its name with the object, and the object is
common.}{A kite, or a bird soaring with wings and tail spread.}

\ufoshapecard{\glampshade}{Lampshade}{Composite; scattered}{A truncated cone,
narrow above and flared below, often reported as glowing from underneath. The
underside glow is the detail that most often turns out to be the light source
itself.}{A streetlamp or floodlight seen through mist, or an illuminated
balloon.}

\ufoshapecard{\glens}{Lenticular}{Discoid; McMinnville, 1950}{Symmetrical,
thicker at the centre, tapering to an edge --- which is precisely what a cap
cloud looks like when it forms over a ridge. The name of the meteorological
phenomenon and the name of the shape class are the same word for a
reason.}{Altocumulus lenticularis, sun through thin cloud, or a lens
artefact.}

\ufoshapecard{\gorb}{Luminous Orb}{Spherical; ubiquitous, all decades}{The
single commonest report in every national archive. A point of light with
apparent volume, no structure, and no scale. It is also the class with the
highest proportion of cases closed by simply asking what was in the sky that
night.}{Venus, Jupiter, a sky lantern, a drone, or a distant landing light.}

\ufoshapecard{\gmanta}{Manta}{Organic; modern reports}{A broad winged form with
a swept trailing edge, reported as undulating rather than banking. Undulation
is the reported feature that most reliably indicates a living thing.}{A large
bird gliding, or a flying-wing airframe at distance.}

\ufoshapecard{\gsphere}{Metallic Sphere}{Spherical; Coyne helicopter case,
1973}{A featureless reflective ball, frequently described as \emph{seamless}
and as showing no exhaust. The absence of features is the whole problem: a
sphere gives you nothing to measure, no aspect angle, and therefore no
distance.}{A high-altitude research or weather balloon, which is spherical,
reflective and silent by design.}

\ufoshapecard{\gslab}{Monolith}{Geometric; Utah and copycats, 2020}{A tall
rectangular slab, usually reported standing rather than flying. The 2020
installations were art, they were admitted as art, and they still generated
sincere anomaly reports for weeks afterwards.}{Sculpture, a survey marker, or
a communications cabinet in poor light.}

\ufoshapecard{\gmushroom}{Mushroom}{Composite; scattered, post-1945}{A domed
cap over a narrower stem. The date matters: the shape enters the record
immediately after the mushroom cloud enters public consciousness, which is a
pattern you will see again and again in this list.}{Cloud formation, a water
tower, or a lit structure with the base in shadow.}

\ufoshapecard{\goct}{Octahedron}{Geometric; scattered modern reports}{Faceted,
angular, described as tumbling or rotating with a repeating glint. The
repeating glint is the giveaway and also the only measurable thing in the
report: time the period.}{A tumbling satellite or rocket body catching the
sun.}

\ufoshapecard{\gpillar}{Pillar of Light}{Luminous; cold-climate reports}{A
vertical column of light with no visible source at either end, reported
standing still for minutes. It is spectacular, it photographs well, and it has
a completely settled meteorological explanation.}{A light pillar --- ground
lighting reflected from ice crystals in still, cold air.}

\ufoshapecard{\gplasma}{Plasma Ball}{Spherical; Hessdalen, from 1981}{A
glowing sphere with a visible corona or flicker, reported in geologically
active valleys. Hessdalen is the one place where instruments were taken to the
phenomenon rather than the other way round, which is why it is worth
respecting.}{Ionised gas from piezoelectric or tectonic stress; ball lightning
remains poorly characterised.}

\ufoshapecard{\gpoint}{Point Source}{Luminous; the base case}{No shape at all:
a light with no resolvable extent, no structure, and no way to establish
distance. It is the honest terminus of any shape taxonomy and, by volume, the
largest single class in every archive there is.}{Aircraft, satellites,
planets, or drones --- an unresolved light, and nothing more.}

\ufoshapecard{\gpyramid}{Pyramid}{Geometric; Russell Island and naval
footage}{A triangular-faced solid with a defined apex, more often reported
stationary than moving. It photographs badly and reads triangular from almost
every angle, which flatters the drawing.}{Defocused point lights --- a
triangular bokeh from a camera iris will do it.}

\ufoshapecard{\gring}{Ring}{Annular; \emph{ovnis en anneau}, France 1970s}{A
thin circle or halo with an apparently empty centre. French investigators
filed a good number of these, and the honest reading is that a ring is what
your eye makes of a defocused bright point.}{A camera artefact, a lit
helicopter rotor arc, or a dispersing contrail.}

\ufoshapecard{\grocket}{Rocket Form}{Cylindrical; twilight launch
windows}{Pointed body with fins and, crucially, a visible plume. This is the
one class that resolves almost perfectly: check the launch schedules for the
hour and the case is normally closed before dinner.}{An actual rocket launch
seen at range, with the plume illuminated by high sunlight.}

\ufoshapecard{\grod}{Rod}{Elongated; camera-artefact era, 1990s}{A long
segmented streak, found in photographs and video and almost never seen by eye.
José Escamilla built a career on these and the answer was available from the
beginning to anyone who understood shutter speed.}{Motion blur on a nearby
insect --- a rolling-shutter and exposure artefact.}

\ufoshapecard{\gsaturn}{Saturn Form}{Discoid; Trindade Island, 1958}{A central
body with an encircling rim or ring. The Trindade photographs made the type
famous and have since been picked apart by everyone who has handled the
negatives. Note that the shape is easy to build and easy to throw.}{A thrown
or suspended model, or a bright object smeared by camera shake.}

\ufoshapecard{\ginsect}{Segmented Body}{Organic; close-range reports}{A jointed
body with limbs or projections, reported only at close range and in poor
light. Close range should mean better data; here it reliably means worse,
because it means one startled witness.}{An insect or bird near the observer,
with scale badly misjudged.}

\ufoshapecard{\gsheet}{Sheet}{Formless; sparse}{A broad flat expanse, reported
as rippling or flexing while it moves. Rigid objects do not ripple, so the
report is either of something flexible or of something being misjudged.}{A
sheet of plastic or tarpaulin lifted on a thermal, or aurora at low
latitude.}

\ufoshapecard{\gsombrero}{Sombrero}{Discoid; Latin American reports}{A wide
flat brim with a raised crown --- the disc family's regional variant, and named
for a hat again. Note how often the catalogue names a shape after whatever the
witness had to hand.}{A lenticular cloud, or a disc-shaped object photographed
from near its plane.}

\ufoshapecard{\gspindle}{Spindle}{Elongated; high-altitude reports}{Tapered to
a point at both ends and reported at great height and speed. The height claim
is the weak link: without a second observer or a radar track, altitude is a
guess dressed as data.}{A high-altitude balloon seen end-on, or a contrail
lit at both ends.}

\ufoshapecard{\gtop}{Spinning Top}{Composite; Brazilian and Japanese
reports}{A conical body with a spindle above, described as rotating on its
axis. It is a toy shape, and I mean that literally: several photographic cases
in this class were resolved to actual toys.}{A thrown or suspended model, or a
lit spinning object close to the camera.}

\ufoshapecard{\gwheel}{Spoked Wheel}{Annular; \emph{Ashtar} wheel, Gulf of
Oman}{A rotating disc with radial spokes of light, reported almost exclusively
at sea. It is one of my favourite entries because it has a beautiful
explanation that most people find less interesting than the mystery.}{Marine
bioluminescence in an internal wave field, seen from the deck.}

\ufoshapecard{\gsquare}{Square}{Geometric; sparse, all decades}{A flat
quadrilateral, reported edge-on as a line and face-on as a panel. It is rare,
and rarity here is a clue: genuinely novel geometry should not be rare if it
is flying about.}{A kite, a lifted tarpaulin or advertising panel, or a
solar-panel glint.}

\ufoshapecard{\gstar}{Star Form}{Luminous; contemporary reports}{A radiating
point with distinct arms, reported as brilliant and often as colour-changing.
This is the class where astronomy closes the file fastest: an app, a star chart
and a timestamp will normally do it.}{Sirius or another bright star scintillating
low in the atmosphere.}

\ufoshapecard{\gteardrop}{Teardrop}{Ovoid; sparse, all decades}{Rounded at one
end, drawn to a point at the other, usually reported as descending. It is the
shape a falling luminous object appears to take because the eye adds a tail to
motion it cannot resolve.}{A meteor or bolide, or burning debris on a ballistic
path.}

\ufoshapecard{\gtee}{Tee}{Angular; sparse}{A crossbar over a stem, typically
reported near airfields and installations. When a shape is reported almost
exclusively where aircraft manoeuvre, treat the location as part of the
evidence.}{An aircraft in a banking turn, or a crane or mast against the sky.}

\ufoshapecard{\gtictac}{Tic Tac}{Cylindrical; USS Nimitz FLIR, 2004}{A smooth
rounded oblong with no wings, no rotor and no plume. This is the modern
flagship case, and my position on it is unchanged: the video is genuine, the
interpretation is not settled, and the sensor is not a witness.}{A distant
aircraft with parallax and glare effects on an infrared sensor.}

\ufoshapecard{\gtophat}{Top Hat}{Composite; European reports, 1954 wave}{A
cylinder set on a flat brim. The 1954 French wave produced a great many of
these, along with the best-documented mass panic in the field's history and
very little in the way of physical evidence.}{An illuminated water tower or
silo, or a lantern seen against low cloud.}

\ufoshapecard{\gtoroid}{Toroid}{Annular; smoke-ring and vortex reports}{A ring
with visible thickness, sometimes reported as rotating about its own axis. It
is a shape that fluid dynamics produces cheaply, which is exactly why it needs
no exotic explanation.}{A vortex ring from an explosion or industrial vent, or
a volcanic gas ring.}

\ufoshapecard{\gtorpedo}{Torpedo}{Cylindrical; naval and coastal sightings}{A
rounded nose tapering to a blunt tail, usually reported low over water and
often as transmedium. The transmedium claim is the interesting part and the
part never captured by more than one sensor at a time.}{A submarine masthead,
a whale surfacing, or an aircraft seen from an unusual angle.}

\ufoshapecard{\gclear}{Transparent Craft}{Formless; rare, modern}{Reported as a
visible distortion or refraction rather than a solid --- ``like heat haze with
edges''. Where a claim is unfalsifiable by construction, it is not evidence,
whatever else it is.}{Heat shimmer, a window or windscreen reflection, or a
soap-film artefact on the lens.}

\ufoshapecard{\gtrapezoid}{Trapezoid}{Geometric; sparse}{A flat form with
unequal parallel edges, reported as tilting. In practice it is nearly always a
rectangle seen in perspective, which is a thing witnesses are demonstrably bad
at correcting for.}{A kite, banner, or panel; or a rectangular object seen
obliquely.}

\ufoshapecard{\gtriangle}{Triangle}{Angular; Belgian wave, 1989--90}{A
flat-bottomed triangle with a light at each vertex, silent, slow, very large.
The Belgian wave produced hundreds of consistent reports and no photograph
anybody can work with, which is the recurring pattern of this field.}{Formation
flying, a stealth or experimental airframe, or three lights read as one
object.}

\ufoshapecard{\gvform}{V Formation}{Formation; Lubbock Lights, 1951}{Lights in
a V, reported as a rigid object rather than a group. Lubbock was resolved to
birds reflecting new street lighting, and the resolution was published, and
almost nobody remembers it.}{Migrating birds catching artificial light from
below.}

\ufoshapecard{\gwedge}{Wedge}{Angular; testing-corridor sightings}{Asymmetric,
thick at one end and knife-edged at the other. Sightings cluster tightly
around known flight-test ranges, and the reports get more detailed the closer
the witness is to a fence they cannot see over.}{A prototype or classified
airframe under test, which is not the same claim as a spacecraft.}

\ufoshapecard{\gwye}{Wye}{Angular; South American reports}{Three arms from a
common hub, described as slowly rotating. It is the shape you would expect if
three lights in loose formation were read as one rigid object, which is
normally what has happened.}{Three lights in formation, or a rotating
navigation beacon.}

\end{multicols}

\rosternarrow
\earthfieldoff
